% Abstract, conclusion, acknowledgment (with the disclosure of AI use), and data availability of Fleet-DT v2.
% RESULTS markers: replace with the numbers of Section VII after the evaluation campaign.

\begin{abstract}
A fleet of unmanned surface vehicles (USVs) shares one Wi-Fi channel with a ground station (GS), and a vessel that sends more data than the GS needs delays the data of the other vessels. This paper presents Fleet-DT, a network digital twin (NDT) that allocates this channel by dynamic bandwidth allocation in steps of 125~ms above the IEEE~802.11 MAC. In each step, each vessel reports the demand of its data and the bytes that its MQTT bridge still holds; the application space of the NDT predicts the capacity of the channel from the busy time that the access point measures and from the efficiency of the distributed coordination function, corrects the prediction with adaptive conformal inference, and admits the data of every vessel within one airtime budget in the order of need of the GS; and the regulator in the MQTT broker of each vessel releases the admitted data at the pace of the radio. The paper proves that the admission returns the lexicographically largest feasible set, that the GS receives every admitted message within the step when the work of the step fits in the realized capacity, and that the long-run fraction of the steps with a missed admitted message exceeds a target by a term that vanishes with the number of steps and a term that counts the steps in which the channel does not carry the required data alone, for every realization of the channel. An evaluation with the unchanged software on emulated hosts over a real-time ns-3 emulation of the channel, with hypotheses and thresholds fixed before the runs, shows RESULTS.
\end{abstract}

\section{Conclusion}\label{sec:conclusion}
This paper presented Fleet-DT, an NDT that allocates the shared Wi-Fi channel of a fleet of USVs among the vessels in steps of 125~ms. In each step, the DTI of Fleet-DT holds the demands, the held bytes, and the measurements of the channel; the AS predicts the capacity of the data window with the efficiency of the DCF and the busy time of the channel, sets the airtime budget with adaptive conformal inference, and admits the units of data of every vessel in the order of need of the GS; and the vessel regulator releases the admitted units at the pace of the radio. The admission returns the lexicographically largest feasible set (Lemma~\ref{lem:greedy}), every admitted message reaches the GS within its step when the work of the step fits in the realized capacity (Lemma~\ref{lem:delivery}), and the long-run fraction of the steps with a missed admitted message stays within the target plus a vanishing term and the steps in which the channel does not carry the required set alone (Theorem~\ref{thm:aci}). RESULTS.

The evaluation covers fleets of up to four vessels in real time; a fleet of more vessels needs an emulation host faster than the host of Section~\ref{sec:method} or an emulation distributed over several hosts. Future work covers a separate transport for the video and steps whose duration follows the demands of the vessels.

\section*{Acknowledgment}
The authors used Claude (Anthropic, model claude-opus-5-5, Cowork) in the preparation of this work: to simulate a review of the conference version and of a draft of this paper against the criteria of IEEE Access; to draft and revise the text of Sections II to VI under the direction of the authors; to verify the proofs of Section IV; to write the software of the emulator, of the vessel regulator, and of the application space, and the scripts of the analysis; and to verify the bibliographic data of the references. The authors defined the problem, the model, and the experiments, reviewed every section, every proof, and every line of code that the paper reports, and take responsibility for the content of the paper.

\section*{Data Availability}
The software of Fleet-DT, the emulator, the frozen inputs of the evaluation (hypotheses, thresholds, seeds, calibration, and analysis scripts), and the data of every run (DTA, logs of the vessels, records of the monitor, ground truth of the emulator, and clock offsets) are available at DOI. The frame-size traces derive from the Pohang Canal Dataset [chung2023pohang] under its CC BY-NC 4.0 license.
